\subsection{Elementary Divisors}

\begin{theorem} \label{thm:operator-primary-cyclic-decomposition} \hlt{Primary Cyclic Decomposition for an Operator}\\
Let $\mu_T=\prod_{i=1}^sp_i^{e_i}$ with $p_i\in\mathbb{F}[z]$ distinct monic irreducible, and set $\mathcal{V}_{p_i}\defeq\ker p_i(T)^{e_i}$. Then
\begin{enumerate}[label=(\roman*)]
\setlength{\itemsep}{0pt}
\item $\mathcal{V}_T=\bigoplus_{i=1}^s\mathcal{V}_{p_i}$;
\item $\mathcal{V}_{p_i}=\bigoplus_{j=1}^{k_i}\mathcal{C}_T(\mathbf{v}_{i,j})$, with $m_{\mathbf{v}_{i,j}}=p_i^{e_{i,j}}$ and $e_i=e_{i,1}\geq e_{i,2}\geq\cdots\geq e_{i,k_i}$;
\item $\operatorname{ElemDiv}(T)=\{p_i^{e_{i,j}}:1\leq i\leq s,\ 1\leq j\leq k_i\}$, corresponding cyclic bases form elementary divisor basis.
\end{enumerate}
\begin{proof}
\begin{enumerate}[label=(\roman*)]
\setlength{\itemsep}{0pt}
\item By Theorem \ref{thm:primary-cyclic-decomposition-pid}, $\mathcal{V}_T=\bigoplus_i\mathcal{V}_T[p_i^\infty]$; $\mathcal{V}_T[p_i^\infty]\defeq\{\mathbf{x}\in\mathcal{V}:\exists N\geq1,\ p_i(T)^N\mathbf{x}=0\}=\ker p_i(T)^{e_i}=\mathcal{V}_{p_i}$.
\item By Theorem \ref{thm:primary-cyclic-decomposition-pid}, $\mathcal{V}_{p_i}=\bigoplus_{j=1}^{k_i}\mathcal{C}_T(\mathbf{v}_{i,j})$, $m_{\mathbf{v}_{i,j}}=p_i^{e_{i,j}}$ and $e_i=e_{i,1}\geq\cdots\geq e_{i,k_i}$.
\item Definition \ref{def:operator-invariant-factors-elementary-divisors} gives $\operatorname{ElemDiv}(T)=\{p_i^{e_{i,j}}:1\leq i\leq s,\ 1\leq j\leq k_i\}$, and concatenating the bases of $\mathcal{C}_T(\mathbf{v}_{i,j})$ gives an elementary divisor basis.
\end{enumerate}
\hlt{Structural Commentary}: Primary module pieces are polynomial-kernel pieces for the operator.\\
\hlt{Key Insight}: Elementary divisors are exactly the orders of the cyclic primary summands.
\end{proof}
\end{theorem}

\begin{theorem} \label{thm:reducing-decompositions-combine-elementary-divisors} \hlt{Reducing Decompositions Combine the Data}\\
Let $\mathcal{V}=\mathcal{A}\oplus\mathcal{B}$ with $\mathcal{A},\mathcal{B}$ $T$-invariant, and set $T_{\mathcal{A}}\defeq T|_{\mathcal{A}}$, $T_{\mathcal{B}}\defeq T|_{\mathcal{B}}$.\\
Then $\mu_T=\operatorname{lcm}(\mu_{T_{\mathcal{A}}},\mu_{T_{\mathcal{B}}})$ and $\operatorname{ElemDiv}(T)=\operatorname{ElemDiv}(T_{\mathcal{A}})\sqcup\operatorname{ElemDiv}(T_{\mathcal{B}})$, where $\sqcup$ is multiset union.\\
The invariant factors are obtained by recombining the elementary divisors into a divisibility chain.
\begin{proof}
By Theorem \ref{thm:operator-submodules-invariant-subspaces}, $\mathcal{V}_T=\mathcal{A}_{T_{\mathcal{A}}}\oplus\mathcal{B}_{T_{\mathcal{B}}}$ as $\mathbb{F}[z]$-modules.\\
Hence $\operatorname{Ann}(\mathcal{V}_T)=\operatorname{Ann}(\mathcal{A}_{T_{\mathcal{A}}})\cap\operatorname{Ann}(\mathcal{B}_{T_{\mathcal{B}}})=\langle\operatorname{lcm}(\mu_{T_{\mathcal{A}}},\mu_{T_{\mathcal{B}}})\rangle$.\\
Proposition \ref{prop:elementary-divisors-direct-sums} gives elementary-divisor identity.\\
\hlt{Structural Commentary}: Reducing decompositions split the operator module.\\
\hlt{Key Insight}: Direct sums take lcm on global orders and multiset union on elementary divisors.
\end{proof}
\end{theorem}

\begin{definition} \label{def:characteristic-polynomial-from-elementary-divisors} \hlt{Characteristic Polynomial from Elementary Divisors}\\
For $T\in\LL(\mathcal{V})$, define $\chi_T(z)\defeq\prod\limits_{q\in\operatorname{ElemDiv}(T)}q(z)$, with product over the multiset $\operatorname{ElemDiv}(T)$.\\
For $A\in\operatorname{Mat}(n,n,\mathbb{F})$, set $\chi_A\defeq\chi_{T_A}$.
\end{definition}

\begin{theorem} \label{thm:minimal-characteristic-elementary-divisor-relations} \hlt{Minimal and Characteristic Polynomials from Elementary Divisors}\\
Let $\operatorname{ElemDiv}(T)=\{p_i^{e_{i,j}}:1\leq i\leq s,\ 1\leq j\leq k_i\}$.\\
Then $\mu_T=\prod_i p_i^{\max_j e_{i,j}}$, $\chi_T=\prod_{i,j}p_i^{e_{i,j}}$, and $\deg\chi_T=\dim\mathcal{V}$.\\
Hence $\mu_T\mid\chi_T$, and $\mu_T,\chi_T$ have the same irreducible factors.
\begin{proof}
By Theorem \ref{thm:operator-primary-cyclic-decomposition}, $\mathcal{V}_T=\bigoplus_{i,j}\mathcal{C}_T(\mathbf{v}_{i,j})$ with $m_{\mathbf{v}_{i,j}}=p_i^{e_{i,j}}$.\\
Hence $\operatorname{Ann}(\mathcal{V}_T)=\langle\operatorname{lcm}_{i,j}p_i^{e_{i,j}}\rangle=\langle\prod_i p_i^{\max_j e_{i,j}}\rangle$, so $\mu_T=\prod_i p_i^{\max_j e_{i,j}}$. Definition \ref{def:characteristic-polynomial-from-elementary-divisors} gives $\chi_T=\prod_{i,j}p_i^{e_{i,j}}$.\\
By Theorem \ref{thm:cyclic-subspace-quotient-and-companion}, $\dim\mathcal{C}_T(\mathbf{v}_{i,j})=\deg p_i^{e_{i,j}}$; summing gives $\deg\chi_T=\dim\mathcal{V}$.\\
Divisibility and common irreducible factors are immediate.\\
\hlt{Structural Commentary}: $\mu_T$ records maximal exponents; $\chi_T$ records all exponents.\\
\hlt{Key Insight}: The pair $(\mu_T,\chi_T)$ forgets the block distribution.
\end{proof}
\end{theorem}

\begin{definition} \label{def:nonderogatory-operator-module} \hlt{Nonderogatory Operators}\\
$T\in\LL(\mathcal{V})$ is \hlt{nonderogatory} if $\mu_T=\chi_T$, equivalently $\deg\mu_T=\dim\mathcal{V}$.\\
For $A\in\operatorname{Mat}(n,n,\mathbb{F})$, $A$ is nonderogatory if $T_A$ is.
\end{definition}

\begin{theorem} \label{thm:nonderogatory-cyclic-equivalences} \hlt{Cyclic and Nonderogatory Are the Same}\\
For $T\in\LL(\mathcal{V})$, the following are equivalent:
\begin{enumerate}[label=(\roman*)]
\setlength{\itemsep}{0pt}
\item $\mathcal{V}_T$ is cyclic;
\item $T$ has a cyclic vector;
\item $\operatorname{ElemDiv}(T)=\{p_1^{e_1},\ldots,p_s^{e_s}\}$, one elementary divisor for each irreducible factor of $\mu_T$;
\item $T$ is nonderogatory.
\end{enumerate}
\begin{proof}
\begin{enumerate}[label=(\roman*).]
\setlength{\itemsep}{0pt}
\item[$(i)$]$\Rightarrow (ii)$ If $\mathcal{V}_T=\mathbb{F}[z]\mathbf{v}$, then $\mathcal{V}=\{p(T)\mathbf{v}:p\in\mathbb{F}[z]\}=\mathcal{C}_T(\mathbf{v})$, so $\mathbf{v}$ is cyclic.
\item[$(ii)$]$\Rightarrow (i)$ If $\mathcal{C}_T(\mathbf{v})=\mathcal{V}$, then $\mathcal{V}_T=\mathbb{F}[z]\mathbf{v}$.
\item[$(i)$]$\Leftrightarrow (iii)$ Theorem \ref{thm:cyclic-torsion-module-characterization} on $\mathcal{V}_T$. Cyclic torsion module has one cyclic summand for each primary factor.
\item[$(iii)$]$\Rightarrow (iv)$ If $\operatorname{ElemDiv}(T)=\{p_i^{e_i}\}_{i=1}^s$, then $\chi_T=\prod_i p_i^{e_i}=\mu_T$, so $T$ is nonderogatory.
\item[$(iv)$]$\Rightarrow (iii)$ Let $\operatorname{ElemDiv}(T)=\{p_i^{e_{i,j}}\}$. If $\mu_T=\chi_T$, then $\prod_i p_i^{\max_j e_{i,j}}=\prod_{i,j}p_i^{e_{i,j}}$, so $k_i=1$ for all $i$.
\end{enumerate}
\hlt{Structural Commentary}: Nonderogatory means all primary information can be generated by one vector.\\
\hlt{Key Insight}: Equality $\mu_T=\chi_T$ says that product and lcm of the cyclic orders coincide.
\end{proof}
\end{theorem}

\begin{perspective} \label{per:elementary-divisors-and-jordan-blocks} \hlt{Elementary Divisors and Jordan Blocks}\\
Over $\mathbb{C}$, elementary divisors are exactly $(z-\lambda)^k$, and $\mathbb{C}[z]/\langle(z-\lambda)^k\rangle$ corresponds to the Jordan block $J_k(\lambda)$.\\
Thus $\operatorname{ElemDiv}(T)$ is the multiset of Jordan blocks.
\end{perspective}
